\documentclass[12pt]{article}
\usepackage{mathtools}
\usepackage{amssymb}
\usepackage{fontspec}
\setmainfont{EB Garamond}

% Set page size and margins
\usepackage[a4paper,top=2cm,bottom=1.5cm,left=1cm,right=1cm]{geometry}
\usepackage{fancyhdr}
\pagestyle{fancy}
\setlength{\parindent}{0pt}

% Useful packages
\usepackage{minted}
\usepackage{graphicx}
\graphicspath{{./media/}}
\usepackage{subfig}
\usepackage{float}
\usepackage[colorlinks=true, allcolors=blue]{hyperref}
\usepackage{siunitx}


\title{\vspace{-2cm} Θεμελιώσεις Κρυπτογραφίας \\
        Homework 2}
\author{Μαρινόπουλος Χριστόφορος 10522}
\date{Εαρινό εξάμηνο 2025}

\lhead{Homework 2}
\chead{Θεμελιώσεις Κρυπτογραφίας}
\rhead{Μαρινόπουλος Χριστόφορος}

\begin{document}
\maketitle
\section*{Άσκηση 1}

Έστω ότι έχουμε ένα προβλέψιμο νόμισμα. Π.χ. η πιθανότητα να έρθει γράμμα είναι \( \frac{1}{2} + \delta \) για μη μηδενικό \(\delta\) με \(\delta \in \left( -\frac{1}{2}, \frac{1}{2} \right)\).
Μπορείτε να βρείτε ένα πείραμα ώστε να μετατρέψετε το προβλέψιμο νόμισμα σε τυχαίο νόμισμα, δηλαδή η πιθανότητα να έρθει γράμμα ή κορώνα να είναι ίδιες;
\section*{Απάντηση}
Για να μετατραπεί το προβλέψιμο νόμισμα σε δίκαιο αρκεί να βρούμε δύο ισοπίθανα γεγονότα και να αντιστοιχίσουμε το ένα στο αποτέλεσμα \textbf{Κορώνα} και το άλλο στο αποτέλεσμα \textbf{Γράμματα}. Τα ισοπίθανα αυτά γεγονότα είναι τα εξής:
\[
\textbf{Pr[ΚΓ]} = \text{Pr[K]} * \text{Pr[Γ]} = \textbf{Pr[ΓΚ]}
\]
Επομένως, αρκεί να πραγματοποιήσουμε το ακόλουθο πείραμα:
\begin{itemize}
    \item Ρίχνουμε το νόμισμα δύο φορές.
    \item \textbf{Εάν το αποτέλεσμα είναι (Κορώνα, Γράμματα)}, καταγράφουμε \textbf{Κορώνα}.
    \item \textbf{Εάν το αποτέλεσμα είναι (Γράμματα, Κορώνα)}, καταγράφουμε \textbf{Γράμματα}.
    \item \textbf{Εάν το αποτέλεσμα είναι (Κορώνα, Κορώνα) ή (Γράμματα, Γράμματα)}, αγνοούμε το ζεύγος και επαναλαμβάνουμε το πείραμα.
\end{itemize}
\hfill \blacksquare
\section*{Άσκηση 2}
Υπό ποιες προϋποθέσεις θα μπορούσε το σύστημα του Καίσαρα να έχει
τέλεια ασφάλεια;
\section*{Απάντηση}
Το σύστημα του Καίσαρα έχει δύο μεγάλα μειονεκτήματα:
\begin{enumerate}
    \item \textbf{Το κάθε γράμμα κρυπτογραφείται με το ίδιο κλειδί} με αποτέλεσμα κάποιος που γνωρίζει το κρυπτογραφημένο μήνυμα και ένα γράμμα του μη κρυπτογραφημένου μηνύματος να έχει τη δυνατότητα να βρει το κλειδί άρα και το μη κρυπτογραφημένο μήνυμα.
    \item \textbf{Το σύστημα κρυπτογραφεί μια λέξη/πρόταση} το οποίο μπορεί να επιτρέψει σε κάποιον που κατέχει ένα μέρος του μη κρυπτογραφημένου μηνύματος, με στατιστική ανάλυση να αποκτήσει περισσότερη πληροφορία για το μη κρυπτογραφημένο μήνυμα.   
\end{enumerate}
Συνεπώς, θα μπορούσε να έχει τέλεια ασφάλεια θα αν:
\begin{enumerate}
    \item \textbf{Το κάθε γράμμα του μηνύματος κρυπτογραφείται με διαφορετικό κλειδί} 
    \item \textbf{Το κρυπτογραφημένο μήνυμα δεν αντιστοιχίζεται άμεσα σε λέξεις ή προτάσεις}
\end{enumerate}
\hfill \blacksquare
\section*{Άσκηση 3}
Έχουμε το εξής σύστημα. Το σύνολο μηνυμάτων και κρυπτογραφημένων
μηνυμάτων είναι {0,1,2,…,100}. To σύνολο των κλειδιών είναι οι (περιττοί) πρώτοι
αριθμοί μικρότεροι του 100. Η συνάρτηση κρυπτογράφησης είναι
όπου p το κλειδί (δηλ. ένας περιττός πρώτος <100).
Να εξετάσετε αν το σύστημα έχει σημασιολογική ασφάλεια.
\section*{Απάντηση}
Σύμφωνα με τον 2ο ορισμό της σημασιολογικής ασφάλειας, ένα κρυπτοσύστημα δεν θεωρείται σημασιολογικά ασφαλές αν, για δύο γνωστά μηνύματα \( m_0 \) και \( m_1\) ένας επιτιθέμενος μπορεί να καθορίσει ποιο μήνυμα από τα δύο κρυπτογραφήθηκε στο \( c \) με πιθανότητα \(>\frac{1}{2}\). 
\newline
\newline
Στο συγκεκριμένο κρυπτοσύστημα για \(m_0 = 0\) και \(m_1 = 1\) για οποιοδήποτε κλειδί \( p \) ο επιτιθέμενος γνωρίζει με πιθανότητα \(1> \frac{1}{2}\) ποιο μήνυμα κρυπτογραφήθηκε. Αυτό ισχύει διότι:
\begin{equation*}
\begin{array}{c}
   0 \mod p = 0\\
   1 \mod p = 1
\end{array}
\end{equation*}
Επομένως, το κρυπτοσύστημα δεν είναι σημασιολογικά ασφαλές.
\hfill \blacksquare
\end{document}